\chapter{NKFont : polices, atlas et métriques}

Afficher du texte est l'opération la plus banale d'une interface et la plus mal
comprise. On croit qu'il suffit d'« écrire une chaîne » ; en réalité, il faut
convertir des octets UTF-8 en \emph{codepoints}, trouver le dessin de chaque
caractère, le rastériser une fois pour toutes dans une grande image, puis
émettre deux triangles texturés par lettre en faisant avancer un curseur.

\texttt{NKFont} fait tout cela, et \textbf{rien d'autre}. C'est un module
remarquablement honnête sur son périmètre : il produit un bitmap en mémoire
vive, et vous rend pour chaque caractère un rectangle à l'écran et un rectangle
dans le bitmap. Ce qu'on en fait ensuite ne le regarde pas.

\section{Le module ne connaît ni image, ni GPU, ni fenêtre}

\begin{nkcode}{Kernel/Runtime/NKFont/NKFont.jenga:35-39}
    nkentseudependson(
        ["NKPlatform", "NKCore", "NKMemory", "NKMath", "NKContainers", "NKThreading", "NKLogger"],
        selfexport="NKFont",
        extra_includes=["src"],
    )
\end{nkcode}

Ni \texttt{NKImage}, ni \texttt{NKCanvas}, ni \texttt{NKWindow}. NKFont est même
plus autonome que NKImage : il ne dépend pas du système de fichiers pour
fonctionner, puisqu'il embarque ses propres polices — nous y venons.

\begin{nkcode}{Kernel/Runtime/NKFont/ — arborescence}
NKFont/
  NKFont.jenga
  src/NKFont/
    NkFont.h            ( 366 l.)  <- EN-TETE PUBLIC PRINCIPAL (NkFontAtlas + NkFont)
    NkFontAtlas.cpp     (~ 700 l.) <- packing + rasterisation
    NkFontMesh.cpp      (~ 400 l.) <- extrusion 3D des glyphes
    Core/
      NkFontTypes.h                <- alias nkft_*, NkFontCodepoint, NkGlyphId
      NkFontParser.h/.cpp          <- parser TTF/OTF/CFF/WOFF + rasteriseur + SDF
      NkFontRasterizer.cpp
      NkFontDetect.h/.cpp          <- NkFontProfile / NkFontDetector
      NkFontSizeCache.h/.cpp       <- cache multi-tailles
    Embedded/
      NkFontEmbedded.h/.cpp        <- 10 polices compressees DANS le binaire
\end{nkcode}

\begin{nkpiege}[La description du \texttt{.jenga} promet six fonctionnalités inexistantes]
\nkfile{Kernel/Runtime/NKFont/NKFont.jenga:7-23} annonce fièrement le support de
BDF, de Type 1 (PFB/PFA), de WOFF2 avec Brotli, d'une machine virtuelle de
\emph{hinting} TrueType, de GSUB pour les ligatures, et de l'algorithme
bidirectionnel UAX\#9.

Une recherche exhaustive sur \nkfile{src/} donne \textbf{zéro occurrence} de
\texttt{BDF}, \texttt{Type1}, \texttt{PFB}, \texttt{Brotli}, \texttt{Bidi},
\texttt{GSUB} et \texttt{Hinting}. La seule mention de \texttt{WOFF2} est un
avertissement disant qu'il n'est pas supporté. \texttt{GPOS} apparaît une fois,
pour lire un décalage de table qui n'est ensuite jamais utilisé
(\nkfile{Core/NkFontParser.cpp:993}).

La liste honnête est celle que le parser affiche lui-même :

\begin{nkcode}{Kernel/Runtime/NKFont/src/NKFont/Core/NkFontParser.h:5-14}
// Formats supportés :
//   OK  TTF SimpleGlyph + CompositeGlyph
//   OK  TTC (TrueType Collection) via faceIndex
//   OK  OTF avec table glyf (contours TrueType)
//   OK  OTF/CFF (Type 2 charstrings — interpréteur intégré)
//   OK  WOFF (décompresseur zlib intégré)
//   /!\ WOFF2 (Brotli — non supporté, dépendance externe requise)
//   OK  cmap format 4 (BMP) et format 12 (full Unicode)
//   OK  kern table format 0
//   OK  SDF generation (Signed Distance Field)
\end{nkcode}

C'est déjà beaucoup — un interpréteur CFF Type 2 écrit à la main n'est pas rien.
Mais c'est \emph{ça}, et pas ce que promet le fichier de build.
\end{nkpiege}

\section{Les quatre objets à connaître}

\subsection{\texttt{NkFontAtlas} — le propriétaire de tout}

L'atlas est la seule chose que vous créez vous-même. Il possède la texture, les
configurations, et les polices.

\begin{nkcode}{Kernel/Runtime/NKFont/src/NKFont/NkFont.h:102-133 (champs publics)}
            void *texID = nullptr;
            nkft_int32 texWidth = 0, texHeight = 0;
            nkft_uint8 *texPixels = nullptr;      // ALPHA8, 1 octet/pixel
            bool texReady = false;
            nkft_int32 texDesiredWidth = 0;
            nkft_int32 texGlyphPadding = 2;       ///< Padding recommandé >= 2 pour éviter le bleeding.
            bool sdfMode = false;
            nkft_int32 sdfSpread = 6;             ///< Rayon SDF en pixels (4-8 typique).
            NkFont *fonts[NK_FONT_ATLAS_MAX_FONTS] = {};
            NkFontConfig configs[NK_FONT_ATLAS_MAX_FONTS];
            nkft_uint32 fontCount = 0u;
\end{nkcode}

Ce sont des champs publics, sans accesseurs : on les lit et on les écrit
directement. C'est délibéré et cohérent avec le reste du moteur.

\begin{nkcode}{Kernel/Runtime/NKFont/src/NKFont/NkFont.h:134-157}
NkFont* AddFontFromFile(const char* path, nkft_float32 sizePixels, const NkFontConfig* cfg = nullptr);
NkFont* AddFontFromMemory(const nkft_uint8* data, nkft_size dataSize,
                          nkft_float32 sizePixels, const NkFontConfig* cfg = nullptr);
NkFont* AddFontFromMemoryOwned(nkft_uint8* data, nkft_size dataSize,
                               nkft_float32 sizePixels, const NkFontConfig* cfg = nullptr);

static const nkft_uint32* GetGlyphRangesDefault();
static const nkft_uint32* GetGlyphRangesLatinExtA();
static const nkft_uint32* GetGlyphRangesCyrillic();
static const nkft_uint32* GetGlyphRangesGreek();
static const nkft_uint32* GetGlyphRangesChineseFull();

bool Build();                                            // <- rasterise TOUT

void GetTexDataAsAlpha8 (nkft_uint8** op, nkft_int32* ow, nkft_int32* oh, nkft_int32* ob = nullptr);
void GetTexDataAsRGBA32 (nkft_uint8** op, nkft_int32* ow, nkft_int32* oh, nkft_int32* ob = nullptr);

void ClearTexData();
void Clear();
bool IsBuilt() const;
\end{nkcode}

Comme \texttt{NkImage}, l'atlas n'est pas copiable :
\texttt{NkFontAtlas(const NkFontAtlas\&) = delete;} (\nkfile{NkFont.h:129}).

\subsection{\texttt{NkFont} — une face rastérisée à une taille}

Attention au nom : dans ce module, \texttt{NkFont} n'est pas « une police ».
C'est \textbf{une police à une taille donnée}, déjà rastérisée. Charger Inter en
18 et en 32 pixels donne \emph{deux} \texttt{NkFont}.

\begin{nkcode}{Kernel/Runtime/NKFont/src/NKFont/NkFont.h:211-219}
            NkFontAtlas *containerAtlas = nullptr;    // l'atlas POSSEDE cette NkFont
            NkFontConfig config;
            nkft_float32 fontSize = 0, ascent = 0, descent = 0, lineAdvance = 0, scale = 1;
            NkFontGlyph glyphs[NK_FONT_MAX_GLYPHS];
            nkft_uint32 glyphCount = 0u;
            const NkFontGlyph *fallbackGlyph = nullptr;
            nkft_float32 fallbackAdvanceX = 0;
\end{nkcode}

Le commentaire d'en-tête est catégorique sur la propriété :

\begin{nkcode}{Kernel/Runtime/NKFont/src/NKFont/NkFont.h:88-90}
 NkFont (le struct produit) n'est pas auto-portant : il est detenu par
 son `containerAtlas` qui parse le TTF, fournit les glyphes et la
 texture. NkFont seul ne « se charge » pas.
\end{nkcode}

\begin{nkretenir}
Les \texttt{NkFont*} que vous recevez appartiennent à l'atlas. \textbf{Ne les
détruisez jamais vous-même}, et ne les gardez pas au-delà de la vie de l'atlas.
Détruire l'atlas détruit toutes ses faces.
\end{nkretenir}

\subsection{\texttt{NkFontGlyph} — la structure qui fait tout le travail}

\begin{nkcode}{Kernel/Runtime/NKFont/src/NKFont/NkFont.h:32-38}
    struct NkFontGlyph {
            NkFontCodepoint codepoint = 0;
            nkft_float32 advanceX = 0.f;
            nkft_float32 x0 = 0, y0 = 0, x1 = 0, y1 = 0; ///< Position quad (pixels écran, relative au curseur).
            nkft_float32 u0 = 0, v0 = 0, u1 = 0, v1 = 0; ///< UV dans l'atlas [0..1].
            bool visible = true;
    };
\end{nkcode}

Lisez-la attentivement : c'est la structure la plus importante du chapitre.

\begin{itemize}
  \item \texttt{x0, y0, x1, y1} sont les coins du rectangle à l'écran,
        \textbf{relatifs au curseur} — pas absolus. Vous ajoutez la position du
        curseur, et c'est tout ;
  \item \texttt{u0, v0, u1, v1} sont les coordonnées dans l'atlas, déjà
        normalisées entre 0 et 1 ;
  \item \texttt{advanceX} est de combien avancer le curseur après ce caractère ;
  \item \texttt{visible} vaut \texttt{false} pour l'espace et les caractères qui
        n'ont pas de dessin : on saute l'émission du quad mais on avance quand
        même le curseur.
\end{itemize}

Il n'y a \textbf{rien d'autre à calculer}. Pas de métriques à recomposer, pas de
conversion d'unités de police en pixels : \texttt{Build()} l'a déjà fait.

\subsection{\texttt{NkFontConfig} — les réglages, posés avant l'ajout}

\begin{nkcode}{Kernel/Runtime/NKFont/src/NKFont/NkFont.h:44-62}
    struct NkFontConfig {
            const nkft_uint8 *fontData = nullptr;
            nkft_size fontDataSize = 0u;
            bool fontDataOwned = false;
            nkft_int32 fontIndex = 0;
            nkft_float32 sizePixels = 13.f;
            nkft_int32 oversampleH = 2;
            nkft_int32 oversampleV = 1;
            bool pixelSnapH = false;
            math::NkVec2f glyphOffset = {0.f, 0.f};
            nkft_float32 glyphMinAdvanceX = 0.f;
            nkft_float32 glyphMaxAdvanceX = 1e9f;
            nkft_float32 glyphExtraSpacing = 0.f;
            const nkft_uint32 *glyphRanges = nullptr;
            NkFontCodepoint glyphFallback = 0x003Fu;
            bool mergeMode = false;
            nkft_float32 rasterizerMultiply = 1.f;
            char name[40] = {};
    };
\end{nkcode}

Dix-sept champs, dont deux comptent vraiment au début.

\texttt{glyphRanges} est un tableau de \emph{paires} \texttt{\{début, fin\}}
terminé par un zéro. Si vous le laissez à \texttt{nullptr}, vous obtenez la
plage par défaut. Si vous voulez le latin étendu plus les caractères de
dessin de boîtes, vous écrivez :

\begin{nkcode}{Exemple écrit pour le cours (API : NkFont.h:56, 134)}
static const uint32 kRanges[] = { 0x0020, 0x00FF, 0x2500, 0x257F, 0 };  // Latin-1 + box-drawing
NkFontConfig cfg;
cfg.glyphRanges = kRanges;
NkFont *f = atlas.AddFontFromFile("Resources/Fonts/Inter-Regular.ttf", 20.f, &cfg);
\end{nkcode}

\texttt{mergeMode} est le mécanisme de \emph{repli} : en ajoutant une seconde
police avec \texttt{mergeMode = true}, ses glyphes viennent combler ceux qui
manquent dans la première, dans la même \texttt{NkFont}. C'est ainsi qu'une
police latine peut afficher des idéogrammes ou des émojis. Nous verrons le code
réel plus loin.

\section{Les dix polices embarquées : écrire du texte sans aucun fichier}

C'est la meilleure porte d'entrée du module, et sans doute la fonctionnalité la
plus sous-estimée du moteur.

\begin{nkcode}{Kernel/Runtime/NKFont/src/NKFont/Embedded/NkFontEmbedded.h:49-74}
enum class NkEmbeddedFontId : nkft_uint32 {
    ProggyClean = 0, ProggyTiny = 1, DroidSans = 2, Karla = 3, Roboto = 4,
    Cousine = 5, SourceCodePro = 6, DroidSerif = 7, DejaVuSansMono = 8, Inter = 9,
    Count = 10, Default = ProggyClean,
};
\end{nkcode}

Ces dix polices sont \emph{réellement} présentes : le fichier
\nkfile{Embedded/NkFontEmbedded.cpp} pèse 2,1 Mo de données compressées, et le
registre (\nkfile{NkFontEmbedded.cpp:20265-20370}) fait pointer chaque entrée sur
un tableau non vide. Aucune n'est un espace réservé.

\begin{nkcode}{Kernel/Runtime/NKFont/src/NKFont/Embedded/NkFontEmbedded.h:120-158}
static bool IsAvailable(NkEmbeddedFontId id);
static const NkEmbeddedFontData* GetData(NkEmbeddedFontId id);
static NkFont* AddToAtlas(NkFontAtlas& atlas, NkEmbeddedFontId id,
                          nkft_float32 sizePx = 0.f, const NkFontConfig* cfg = nullptr);
static NkFont* AddDefaultFont(NkFontAtlas& atlas, const NkFontConfig* cfg = nullptr);
static const char* GetName(NkEmbeddedFontId id);
static const NkEmbeddedFontData* GetAll(nkft_int32* outCount);
static nkft_uint8* DecompressData(const NkEmbeddedFontData& data, nkft_uint32* outSize = nullptr);
static void FreeDecompressedData(nkft_uint8* ptr);
\end{nkcode}

Les licences sont déclarées dans le registre : ProggyClean et ProggyTiny sont en
MIT ; DroidSans, Roboto, Cousine et DroidSerif en Apache 2.0 ; Karla,
SourceCodePro et Inter en OFL 1.1 ; DejaVuSansMono en Bitstream Vera / domaine
public. Ce n'est pas un détail décoratif : cela veut dire que vous pouvez
distribuer votre exécutable sans vous poser de question.

\begin{nkretenir}
\textbf{Votre première application affichant du texte n'a besoin d'aucun fichier
de police, d'aucun dossier de ressources, d'aucun chemin relatif.} Une ligne
suffit :
\texttt{NkFontEmbedded::AddToAtlas(atlas, NkEmbeddedFontId::Inter, 18.f)}.
Vous n'aurez jamais l'erreur « police introuvable ».
\end{nkretenir}

Le dépôt s'en sert d'ailleurs avec un repli en cascade, parce que rien n'oblige
une police donnée à être compilée dans un binaire particulier :

\begin{nkcode}{Applications/Pong/src/Pong/Render/FontAtlas.cpp:52-58}
            NkEmbeddedFontId fontId = NkEmbeddedFontId::ProggyClean;
            if (NkFontEmbedded::IsAvailable(NkEmbeddedFontId::Karla)) {
                fontId = NkEmbeddedFontId::Karla;
            } else if (NkFontEmbedded::IsAvailable(NkEmbeddedFontId::DroidSans)) {
                fontId = NkEmbeddedFontId::DroidSans;
            } else if (NkFontEmbedded::IsAvailable(NkEmbeddedFontId::Roboto)) {
                fontId = NkEmbeddedFontId::Roboto;
            }
\end{nkcode}

\section{La séquence canonique en quatre temps}

Tout usage de NKFont suit le même ordre, sans exception :

\begin{enumerate}
  \item \texttt{atlas.Clear()} — si l'atlas a déjà servi ;
  \item \texttt{Add*()} — une ou plusieurs fois, une par police et par taille ;
  \item \texttt{atlas.Build()} — \textbf{c'est ici que les pixels apparaissent} ;
  \item \texttt{atlas.GetTexDataAsAlpha8(...)} — on récupère le bitmap.
\end{enumerate}

Le pont officiel entre NKFont et NKGui l'applique littéralement :

\begin{nkcode}{Kernel/Runtime/NKGui/src/NKGui/Core/NkGuiFont.cpp:118-133}
        bool NkGuiFont::LoadEmbedded(NkEmbeddedFontId id, float32 sizePx, bool extFallback) noexcept {
            atlas.Clear();
            face = nullptr;
            pixels = nullptr; // rechargeable
            NkFontConfig cfg;
            cfg.glyphRanges = NkGlyphRanges();
            face = NkFontEmbedded::AddToAtlas(atlas, id, sizePx, &cfg);
            if (!face)
                return false;
            NkMergeFallback(atlas, sizePx, extFallback); // repli pour les glyphes manquants
            if (!atlas.Build())
                return false;
            int32 bpp = 0;
            atlas.GetTexDataAsAlpha8(&pixels, &atlasW, &atlasH, &bpp);
            dirty = (pixels != nullptr && atlasW > 0 && atlasH > 0);
            return dirty;
        }
\end{nkcode}

\begin{nkpiege}[L'ordre n'est pas négociable, et l'échec est silencieux]
\texttt{Build()} refuse de travailler sur un atlas vide et le journalise :

\begin{nkcode}{Kernel/Runtime/NKFont/src/NKFont/NkFontAtlas.cpp:330-333}
        if (fontCount == 0) {
            logger.Info("[NkFontAtlas] Build():aucune fonte\n");
            return false;
        }
\end{nkcode}

Mais \texttt{GetTexDataAsAlpha8}, appelé avant \texttt{Build()}, ne journalise
\emph{rien} : il pose simplement \texttt{nullptr}
(\nkfile{NkFontAtlas.cpp:640-642} : \texttt{if (!texReady || !texPixels) \{ *op =
nullptr; ... \}}). Vous uploadez alors un pointeur nul, le backend ne fait rien,
et vous cherchez pendant une heure pourquoi tous vos textes sont invisibles.

\textbf{Testez toujours le pointeur récupéré}, comme le fait \texttt{NkGuiFont}
avec son champ \texttt{dirty}.
\end{nkpiege}

\subsection{Le mécanisme de repli en action}

\begin{nkcode}{Kernel/Runtime/NKGui/src/NKGui/Core/NkGuiFont.cpp:98-113}
        static void NkMergeFallback(NkFontAtlas &atlas, float32 sizePx, bool ext) noexcept {
            if (!ext)
                return;
            auto add = [&](const char *path, const uint32 *ranges) {
                if (!NkFileExists(path))
                    return;
                NkFontConfig fb;
                fb.glyphRanges = ranges;
                fb.mergeMode = true;
                atlas.AddFontFromFile(path, sizePx > 0.f ? sizePx : 16.f, &fb);
            };
            add(gFbBroad, NkBroadRanges());
            add(gFbCjk, NkCjkRanges());
            add(gFbEmoji, NkEmojiRanges());
        }
\end{nkcode}

Trois polices de repli : une couverture large, une pour le chinois-japonais-
coréen, une pour les émojis. Chacune n'est ajoutée que si le fichier existe —
d'où le \texttt{NkFileExists}. Les chemins, eux, sont posés par l'application :

\begin{nkcode}{Kernel/Runtime/NKGui/src/NKGui/Core/NkGuiFont.h:76-78}
 A poser par l'APPLICATION (ex. NKCode, depuis son dossier data/fonts) AVANT
 de charger les polices. Un chemin nullptr/vide = role desactive.
void NkSetFallbackFontPaths(const char* broad, const char* cjk, const char* emoji) noexcept;
\end{nkcode}

\emph{Avant} de charger — pas après. Une fois l'atlas construit, il est trop
tard.

\section{Dessiner une chaîne de caractères}

\subsection{Le patron universel}

Voici la boucle que vous écrirez, sous une forme ou une autre, chaque fois que
vous rendrez du texte vous-même. C'est celle de NKGui :

\begin{nkcode}{Kernel/Runtime/NKGui/src/NKGui/Core/NkGuiDrawList.cpp:245-266}
            while (p < end) {
                const NkFontCodepoint cp = NkFontDecodeUTF8(&p, end);
                if (cp == 0u)
                    break;
                const NkFontGlyph *g = face->FindGlyph(cp);
                if (!g)
                    continue;
                if (g->visible) {
                    const float32 x0 = NkGuiPixelSnap(x + g->x0), y0 = NkGuiPixelSnap(y + g->y0);
                    const float32 x1 = x0 + (g->x1 - g->x0), y1 = y0 + (g->y1 - g->y0);
                    if (x1 > xEnd)
                        break; // troncature simple
                    const uint32 i0 = Vtx({x0 + sTop, y0}, {g->u0, g->v0}, c);
                    const uint32 i1 = Vtx({x1 + sTop, y0}, {g->u1, g->v0}, c);
                    const uint32 i2 = Vtx({x1 + sBot, y1}, {g->u1, g->v1}, c);
                    const uint32 i3 = Vtx({x0 + sBot, y1}, {g->u0, g->v1}, c);
                    Tri(i0, i1, i2, texId);
                    Tri(i0, i2, i3, texId);
                }
                x += g->advanceX;
            }
\end{nkcode}

Décortiquons.

\begin{enumerate}
  \item \texttt{NkFontDecodeUTF8(\&p, end)} avance le pointeur \texttt{p} et
        rend le \emph{codepoint}. C'est une fonction libre du module
        (\nkfile{NkFont.h:313}), doublée d'une méthode statique
        \texttt{NkFont::DecodeUTF8} (\nkfile{NkFont.h:235}) ;
  \item \texttt{FindGlyph(cp)} cherche le glyphe \textbf{avec repli} : si le
        caractère n'existe pas dans la police, on obtient le glyphe de
        remplacement plutôt que \texttt{nullptr}. La variante
        \texttt{FindGlyphNoFallback} existe si vous voulez savoir la vérité ;
  \item si \texttt{visible}, on émet quatre sommets et deux triangles. Les
        positions sont \texttt{x + g->x0} et \texttt{y + g->y0} : le curseur plus
        le décalage du glyphe ;
  \item \textbf{en dehors du \texttt{if}}, on avance : \texttt{x += g->advanceX}.
        L'espace n'a pas de dessin mais fait avancer le curseur.
\end{enumerate}

\subsection{Le pixel-snap, et l'invariant subtil}

Regardez à nouveau : les positions passent par \texttt{NkGuiPixelSnap}. Le
commentaire explique pourquoi :

\begin{nkcode}{Kernel/Runtime/NKGui/src/NKGui/Core/NkGuiDrawList.cpp:225-231}
 Le curseur de texte accumule des avances FRACTIONNAIRES. Sans arrondi,
 seul le PREMIER glyphe d'une chaîne tombe sur un pixel entier ; tous les
 suivants dérivent et échantillonnent l'atlas ENTRE deux texels, ce qui
 rend le texte uniformément flou.
\end{nkcode}

Mais l'arrondi est appliqué avec une précaution qu'on ne devine pas :

\begin{nkcode}{Kernel/Runtime/NKGui/src/NKGui/Core/NkGuiDrawList.cpp:257-266}
 On arrondit la POSITION du quad, jamais l'AVANCE : `x` continue
 d'accumuler la valeur exacte, donc la largeur totale de la chaîne
 est inchangée et MeasureWidth reste d'accord avec le rendu.
\end{nkcode}

\begin{nkretenir}[Arrondir l'affichage, pas le calcul]
Si vous arrondissiez \texttt{advanceX}, l'erreur s'accumulerait caractère après
caractère et la largeur mesurée ne correspondrait plus à la largeur dessinée —
vos alignements, vos centrages et vos troncatures deviendraient tous faux. C'est
une leçon générale : \textbf{on arrondit à l'affichage, jamais dans
l'accumulateur}.
\end{nkretenir}

\subsection{La ligne de base, pas le haut du texte}

\texttt{y} dans la boucle ci-dessus n'est pas le haut du texte : c'est la
\textbf{ligne de base} — la ligne sur laquelle « posent » les lettres, et sous
laquelle descendent les jambages du \texttt{p} et du \texttt{g}. Les
\texttt{g->y0} sont donc négatifs pour la plupart des caractères.

Les métriques nécessaires sont dans la \texttt{NkFont}, calculées par
\texttt{Build()} :

\begin{itemize}
  \item \texttt{ascent} — hauteur au-dessus de la ligne de base ;
  \item \texttt{descent} — profondeur en dessous ;
  \item \texttt{lineAdvance} — de combien descendre pour la ligne suivante ;
  \item \texttt{fontSize}, \texttt{scale} — la taille demandée et le facteur
        appliqué.
\end{itemize}

Donc, pour dessiner à partir d'un coin haut-gauche :

\begin{nkcode}{Exemple écrit pour le cours (API : NkFont.h:213)}
float x = posX;
float y = posY + font->ascent;      // y = LIGNE DE BASE, pas le haut !
// ... boucle de glyphes ...
y += font->lineAdvance;             // ligne suivante
\end{nkcode}

\begin{nkpiege}[La constante magique de Pong]
Le jeu \texttt{Pong} approxime la montée par un facteur empirique :

\begin{nkcode}{Applications/Pong/src/Pong/Render/FontAtlas.cpp:206-208}
            const float ascent = fontPx * 0.82f;
\end{nkcode}

Cela fonctionne pour \emph{cette} police, à \emph{ces} tailles. Changez de
police et tout se décale verticalement. La valeur exacte existe pourtant :
c'est \texttt{font->ascent}, calculée par \texttt{Build()}
(\nkfile{NkFontAtlas.cpp:377}). Utilisez-la.
\end{nkpiege}

\subsection{Mesurer avant de dessiner}

\begin{nkcode}{Kernel/Runtime/NKFont/src/NKFont/NkFont.h:229-235}
            nkft_float32 GetCharAdvance(NkFontCodepoint c) const {
                const NkFontGlyph *g = FindGlyph(c);
                return g ? g->advanceX : fallbackAdvanceX;
            }

            nkft_float32 CalcTextSizeX(const char *text, const char *textEnd = nullptr) const;
            static NkFontCodepoint DecodeUTF8(const char **text, const char *textEnd);
\end{nkcode}

\texttt{CalcTextSizeX} parcourt la chaîne et somme les avances. C'est ce qu'il
faut pour centrer un libellé, dimensionner un bouton, ou décider d'une
troncature. Notez qu'il n'existe pas de \texttt{CalcTextSizeY} : la hauteur d'une
ligne, c'est \texttt{lineAdvance}, point.

\section{De l'atlas à l'écran}

\subsection{L'atlas est en alpha 8 bits}

\texttt{GetTexDataAsAlpha8} rend un octet par pixel : la \emph{couverture} du
glyphe, de 0 (rien) à 255 (plein). Ce n'est pas une image en niveaux de gris à
afficher telle quelle — c'est un masque.

Or les rasteriseurs 2D du moteur échantillonnent des textures RGBA. Il faut donc
convertir, et le dépôt explique exactement pourquoi la recette est celle-là :

\begin{nkcode}{Applications/Pong/src/Pong/Render/FontAtlas.cpp:92-104}
            // L'atlas NKFont est alpha8 (1 canal = couverture du glyphe). Le
            // shader 2D de NKCanvas echantillonne la texture en RGBA puis la
            // multiplie par la couleur du vertex. On convertit donc l'atlas en
            // RGBA8 "blanc + alpha=couverture" : texture(1,1,1,cov) * color =
            // (color.rgb, color.a*cov) => exactement le rendu de texte voulu.
            const usize pixCount = static_cast<usize>(w) * static_cast<usize>(h);
            NkVector<uint8> rgba;
            rgba.Resize(pixCount * 4u);
            for (usize i = 0; i < pixCount; ++i) {
                rgba[i * 4 + 0] = 255;
                rgba[i * 4 + 1] = 255;
                rgba[i * 4 + 2] = 255;
                rgba[i * 4 + 3] = pixels[i]; // couverture du glyphe
            }
\end{nkcode}

Blanc partout, alpha = couverture. Multiplié par la couleur du sommet, cela
donne exactement la couleur voulue avec la bonne transparence. Vous n'avez pas à
écrire cette boucle si vous ne le voulez pas : \texttt{GetTexDataAsRGBA32}
(\nkfile{NkFontAtlas.cpp:662-680}) fait rigoureusement la même chose, en allouant
un second tampon à la demande.

\subsection{Le piège d'upload}

\begin{nkcode}{Applications/Pong/src/Pong/Render/FontAtlas.cpp:107-109}
            // NB : surtout PAS Create()+Update() : Update exige gTextureBackend.Update
            // qui n'est pas cable sur tous les backends -> echec silencieux.
\end{nkcode}

Ce commentaire vaut avertissement général : dans NkCanvas, préférez
\texttt{LoadFromImage} en un seul appel plutôt que la paire
\texttt{Create} + \texttt{Update}, tant que tous les backends ne câblent pas
\texttt{Update}.

\subsection{Le chemin NKGui, celui que vous emploierez}

Dans une interface NKGui, tout ce qui précède est encapsulé dans un petit
\emph{wrapper} :

\begin{nkcode}{Kernel/Runtime/NKGui/src/NKGui/Core/NkGuiFont.h:19-69 (abrégé)}
        struct NkGuiFont {
                NkFontAtlas atlas;            ///< possède la texture + les glyphes
                NkFont *face = nullptr;        ///< face produite (détenue par l'atlas)
                uint32 texId = 0x4E4B4654u; ///< 'NKFT' — id stable pour le backend
                uint8 *pixels = nullptr;    ///< atlas alpha8 (détenu par l'atlas)
                int32 atlasW = 0;
                int32 atlasH = 0;
                bool dirty = false; ///< à (ré)uploader côté backend
                bool LoadEmbedded(NkEmbeddedFontId id, float32 sizePx, bool extFallback = true) noexcept;
                bool LoadFromFile(const char *path, float32 sizePx, bool extFallback = true) noexcept;
                /* Face() TexId() Valid() Ascent() Descent() LineHeight() MeasureWidth() */
        };
\end{nkcode}

et son branchement tient en six lignes :

\begin{nkcode}{Applications/NKGuiDemo/src/NKGuiDemo/main.cpp:281-288}
    auto fontPtr = memory::NkMakeUnique<NkGuiFont>();
    if (!fontPtr->LoadEmbedded(NkEmbeddedFontId::DroidSans, 18.f)) {
        fontPtr->LoadEmbedded(NkEmbeddedFontId::ProggyClean, 16.f);
    }
    ctx.font = fontPtr.Get();
    if (fontPtr->Valid()) {
        backend.UploadFontGray8(fontPtr->TexId(), fontPtr->pixels, fontPtr->atlasW, fontPtr->atlasH);
    }
\end{nkcode}

Notez \texttt{UploadFontGray8} et non \texttt{UploadImageRGBA} : le backend sait
qu'un atlas de police est en un seul canal et fait la conversion lui-même
(\nkfile{Kernel/Runtime/NKCanvas/src/NKCanvas/UI/NkGuiCanvasBackend.h:41}).
Après quoi \texttt{ctx.font} pointe la face, et tous les widgets de NKGui
dessinent leur texte tout seuls.

\begin{nkretenir}[Trois lignes, trois responsabilités]
\begin{enumerate}
  \item \texttt{fontPtr->LoadEmbedded(id, taille)} — NKFont rastérise ;
  \item \texttt{backend.UploadFontGray8(...)} — le backend crée la texture GPU ;
  \item \texttt{ctx.font = fontPtr.Get()} — NKGui sait quelle face utiliser.
\end{enumerate}
Oubliez la deuxième : tous les textes sont invisibles, sans message d'erreur.
Oubliez la troisième : les widgets n'affichent aucun libellé, sans message
d'erreur non plus.
\end{nkretenir}

\subsection{Changer d'échelle : recharger et ré-uploader}

Sur un écran à 150\,\%, il ne suffit pas d'agrandir les quads : l'atlas doit être
rastérisé plus gros, sinon le texte est flou. La démo le fait à la touche F9 :

\begin{nkcode}{Applications/NKGuiDemo/src/NKGuiDemo/main.cpp:461-466}
        if (pendingScale > 0.f) { // F9 : appliquer la nouvelle echelle DPI
            SetUiScale(ctx, pendingScale);
            if (fontPtr->LoadEmbedded(NkEmbeddedFontId::DroidSans, 18.f * pendingScale))
                backend.UploadFontGray8(fontPtr->TexId(), fontPtr->pixels, fontPtr->atlasW, fontPtr->atlasH);
            pendingScale = 0.f;
        }
\end{nkcode}

Le ré-upload est obligatoire, et le module le dit :

\begin{nkcode}{Kernel/Runtime/NKFont/src/NKFont/Core/NkFontSizeCache.h:62-63}
 @note Le rebuild de l'atlas invalide toutes les textures GPU.
       L'application doit être notifiée pour re-uploader la texture.
\end{nkcode}

C'est la raison d'être du drapeau \texttt{dirty} de \texttt{NkGuiFont} et des
méthodes \texttt{NeedsGpuUpload()} / \texttt{ClearGpuUploadFlag()} du cache de
tailles.

\section{Le mode SDF : du texte net à toute taille}

Un atlas classique est rastérisé à une taille précise ; l'agrandir le rend flou.
Le mode \emph{Signed Distance Field} stocke, pour chaque texel, la distance au
bord du glyphe plutôt que sa couverture. Un \emph{shader} reconstitue alors un
bord net à n'importe quelle échelle.

\begin{nkcode}{Exemple écrit pour le cours (API : NkFont.h:112-113, 149)}
NkFontAtlas atlas;
atlas.sdfMode   = true;      // AVANT Build()
atlas.sdfSpread = 6;         // 4-8 pixels
NkFontEmbedded::AddToAtlas(atlas, NkEmbeddedFontId::Inter, 48.f);
atlas.Build();
// Utiliser le shader fourni : nkentseu::kNkFontFragSDF (NkFont.h:343)
// avec uSmoothing ~ 0.1 / fontSize.
\end{nkcode}

Ce mode est réellement branché : \nkfile{NkFontAtlas.cpp:496-535} alloue un
tampon temporaire et appelle \texttt{nkfont::NkMakeSDFFromBitmap(...)}.

Et le module fournit les deux \emph{shaders} correspondants, en constantes
compilées :

\begin{nkcode}{Kernel/Runtime/NKFont/src/NKFont/NkFont.h:324-336}
    static constexpr const char *kNkFontFragNormal = R"GLSL(
#version 460 core
in vec2 vUV;
in vec4 vColor;
out vec4 fragColor;
layout(binding=0) uniform sampler2D uAtlas;
void main() {
    float alpha = texture(uAtlas, vUV).a;
    if (alpha < 0.01) discard;
    fragColor = vec4(vColor.rgb, vColor.a * alpha);
}
)GLSL";
\end{nkcode}

\texttt{kNkFontFragSDF} (\nkfile{NkFont.h:343}) est son pendant pour le mode SDF,
avec un uniforme \texttt{uSmoothing}. Vous n'en aurez pas besoin dans une
interface NKGui — le backend a déjà les siens — mais ils sont là si vous écrivez
votre propre rendu.

\section{Les limites dures, et l'invariant du \texttt{Build()}}

Voici la partie que le lecteur pressé saute et regrette.

\subsection{Des tableaux de taille fixe}

\begin{nkcode}{Kernel/Runtime/NKFont/src/NKFont/NkFont.h:99-100, 169-170}
static constexpr nkft_uint32 NK_FONT_ATLAS_MAX_FONTS = 16;
static constexpr nkft_uint32 NK_FONT_ATLAS_MAX_CUSTOM_RECTS = 64;
static constexpr nkft_uint32 NK_FONT_MAX_GLYPHS = 4096u;
static constexpr nkft_uint32 NK_FONT_INDEX_SIZE = 256u;
\end{nkcode}

\texttt{NkFont::glyphs} est un tableau \emph{membre} de 4096 entrées, pas un
conteneur dynamique. Conséquences concrètes :

\begin{itemize}
  \item une face ne peut pas dépasser 4096 glyphes. La plage
        \texttt{GetGlyphRangesChineseFull()} en dépasse largement : elle ne
        rentrera pas ;
  \item un atlas ne peut pas contenir plus de 16 entrées — et une entrée, c'est
        \emph{une police à une taille}. Cinq tailles de deux polices, plus trois
        replis, et vous êtes à treize ;
  \item chaque \texttt{NkFont} est un objet volumineux. Ne les copiez pas ;
        manipulez des \texttt{NkFont*}, ce que l'API impose de toute façon.
\end{itemize}

\subsection{\texttt{Build()} n'est ni réentrant ni parallélisable}

\begin{nkcode}{Kernel/Runtime/NKFont/src/NKFont/NkFontAtlas.cpp:336-337}
        static nkfont::NkFontFaceInfo faceInfos[NK_FONT_ATLAS_MAX_FONTS];
        static nkft_float32 scales[NK_FONT_ATLAS_MAX_FONTS];
\end{nkcode}

Ces tableaux sont \texttt{static} — partagés par toutes les instances. Et pire,
chaque face en garde l'adresse :

\begin{nkcode}{Kernel/Runtime/NKFont/src/NKFont/NkFontAtlas.cpp:381-382}
                // Stockage du faceInfo pour l'extraction des contours
                font->m_FaceInfo = &faceInfos[i];
\end{nkcode}

Les tableaux de \emph{packing} sont eux aussi statiques
(\nkfile{NkFontAtlas.cpp:381-386}).

\begin{nkpiege}[Un seul atlas, un seul \texttt{Build()}, un seul fil]
Trois conséquences, toutes silencieuses :
\begin{enumerate}
  \item deux \texttt{Build()} \textbf{en parallèle} sur deux fils écrivent dans
        les mêmes tableaux : corruption, sans message ;
  \item un \texttt{Build()} sur un atlas B \textbf{invalide} les
        \texttt{m\_FaceInfo} des faces de l'atlas A. Les fonctions qui s'en
        servent — \texttt{GetGlyphOutlinePoints}, l'extrusion 3D — deviennent
        fausses. L'atlas A continue pourtant d'afficher du texte normalement :
        seule l'extraction de contours ment ;
  \item \texttt{Build()} n'est pas réentrant : ne l'appelez pas depuis un
        \emph{callback} déclenché par lui-même.
\end{enumerate}
\textbf{La règle à suivre : un seul \texttt{NkFontAtlas} par application, sur le
fil principal.} C'est ce que fait NKGui, et c'est ce que vous ferez.
\end{nkpiege}

\subsection{Le \texttt{mergeMode} duplique les pointeurs}

\begin{nkcode}{Kernel/Runtime/NKFont/src/NKFont/NkFontAtlas.cpp:185-189}
 Les polices FUSIONNÉES (mergeMode) partagent le MÊME NkFont* que leur police de
 base (cf. AddFontFromMemoryOwned) : fonts[] contient donc des pointeurs DUPLIQUÉS.
 Il faut ne détruire chaque pointeur unique QU'UNE fois — sinon double-free / tas
 corrompu au rechargement (crash zoom).
\end{nkcode}

Vous n'aurez pas à gérer cela vous-même — l'atlas s'en charge — mais retenez le
symptôme : « crash au zoom » signifie « rechargement d'atlas avec fusion ». Si
vous écrivez votre propre wrapper et que vous parcourez \texttt{atlas.fonts[]}
pour faire quelque chose sur chaque face, \textbf{vous visiterez plusieurs fois
le même pointeur}.

\subsection{Le cache de tailles rend \texttt{nullptr} la première fois}

Si vous utilisez \texttt{NkFontSizeCache} (\nkfile{Core/NkFontSizeCache.h:85}),
lisez ce contrat avant :

\begin{nkcode}{Kernel/Runtime/NKFont/src/NKFont/Core/NkFontSizeCache.h:114-117}
 Si la taille n'est pas dans le cache :
   1. Ajoute la fonte à l'atlas.
   2. Marque l'atlas comme nécessitant un rebuild.
   3. Retourne nullptr jusqu'au prochain appel après BuildAtlas().
\end{nkcode}

Un \texttt{GetFont(24.f)} qui renvoie \texttt{nullptr} n'est donc pas forcément
une erreur : c'est peut-être « revenez à la frame suivante ». Le cache plafonne à
\texttt{MAX\_CACHED\_SIZES = 16} tailles (\nkfile{NkFontSizeCache.h:87}).

\section{Ce que le module ne fait pas}

Pour être complet, et pour vous éviter de chercher :

\begin{center}
\begin{tabular}{@{}ll@{}}
\toprule
\textbf{Fonctionnalité} & \textbf{État réel} \\
\midrule
WOFF2 / Brotli        & absent (\nkfile{Core/NkFontParser.h:11}) \\
GSUB, ligatures       & absent \\
Bidi (UAX\#9)         & absent \\
\emph{Hinting} TrueType & absent \\
BDF, Type 1 (PFB/PFA) & absent \\
GPOS                  & offset lu, jamais exploité \\
Kerning par paire     & existe au parser, non remonté \\
Faux-gras             & non — charger une fonte grasse \\
\nkfile{NkFontSdf.h}  & annoncé « futur », n'existe pas \\
\bottomrule
\end{tabular}
\end{center}

Deux précisions.

Le \textbf{kerning} est un vrai cas curieux : le parser sait lire la table
\texttt{kern} format 0 et expose \texttt{NkGetGlyphKernAdvance}
(\nkfile{Core/NkFontParser.h:171}), mais l'information ne remonte pas jusqu'à
\texttt{NkFontGlyph}. Le wrapper NkCanvas l'écrit franchement :
« Le module NKFont n'expose pas le kerning par paire $\rightarrow$ 0 »
(\nkfile{Kernel/Runtime/NKCanvas/src/NKCanvas/Renderer/Resources/NkFont.h:87}).
Un « AV » restera donc plus espacé que dans un traitement de texte.

Le \textbf{gras} n'est pas synthétisé : « bold ignoré (le module ne fait pas de
faux-gras ; charger une fonte bold dédiée si nécessaire) »
(\nkfile{.../Resources/NkFont.h:84-85}). Si vous voulez du gras, ajoutez une
seconde entrée dans l'atlas — en gardant à l'esprit la limite de seize.

\section{L'autre wrapper : \texttt{renderer::NkFont}}

Il existe une seconde façade, côté NkCanvas, d'inspiration SFML :

\begin{nkcode}{Kernel/Runtime/NKCanvas/src/NKCanvas/Renderer/Resources/NkFont.h:63-115 (abrégé)}
        class NkFont {
            public:
                bool LoadFromFile(NkIRenderer2D &renderer, const char *path);
                bool LoadFromMemory(NkIRenderer2D &renderer, const void *data, usize sizeBytes);
                const NkGlyph &GetGlyph(uint32 codepoint, uint32 characterSize, bool bold = false) const;
                float32 GetKerning(uint32 first, uint32 second, uint32 characterSize) const;
                float32 GetLineHeight(uint32 characterSize) const;
                const NkTexture *GetAtlasTexture(uint32 characterSize) const;
                bool IsValid() const;
                void Destroy();
            private:
                struct Page {                       // une page = une TAILLE de caractère
                        uint32 characterSize = 0;
                        ::nkentseu::NkFontAtlas *atlas = nullptr;
                        ::nkentseu::NkFont *moduleFont = nullptr;
                        NkTexture texture;
                };
        };
\end{nkcode}

Attention à l'homonymie : \texttt{renderer::NkFont} n'est pas
\texttt{nkentseu::NkFont}. Le premier est un gestionnaire de pages qui délègue au
second :

\begin{nkcode}{Kernel/Runtime/NKCanvas/src/NKCanvas/Renderer/Resources/NkFont.h:55-62}
 Depuis le refactoring 2026-05-28, NKCanvas ne réimplémente plus la
 rasterisation (ex-FreeType supprimé). renderer::NkFont délègue au module
 NKFont (nkentseu::NkFontAtlas + nkentseu::NkFont) : une "page" (atlas
 rasterisé + texture GPU) par taille de caractère demandée.
\end{nkcode}

Une page par taille demandée : chaque nouvelle \texttt{characterSize} construit
un nouvel atlas. Pratique, mais rappelez-vous l'invariant du \texttt{Build()}
statique — c'est un point à surveiller si vous mélangez ce wrapper avec un atlas
à vous.

\section{Récapitulatif des trois chemins}

\begin{center}
\begin{tabular}{@{}lll@{}}
\toprule
\textbf{Chemin} & \textbf{Quand} & \textbf{Ce que vous écrivez} \\
\midrule
\texttt{NkGuiFont} & interface NKGui & 3 lignes, rien de plus \\
\texttt{renderer::NkFont} + \texttt{NkText} & scène NkCanvas & chargement + dessin \\
Atlas manuel & rendu maison & atlas, conversion, quads \\
\bottomrule
\end{tabular}
\end{center}

Le premier est celui de ce cours. Le troisième est celui qui \emph{explique}, et
c'est pourquoi nous l'avons détaillé : le jour où votre texte est flou, décalé
ou invisible, c'est dans ces quatre étapes que le problème se trouve.

\begin{nkbilan}
NKFont ne connaît ni image, ni GPU, ni fenêtre : il transforme une police en
\emph{glyphes} rangés dans un atlas, et vous rend des coordonnées. Le passage à
l'écran est votre affaire --- ou celle du pont. La séquence est toujours la
même : charger, demander les caractères voulus, construire l'atlas, dessiner ; et
le \texttt{Build()} porte un invariant qu'on ne peut pas contourner --- ce qui
n'a pas été demandé avant ne sera pas dans l'atlas. Le mode \textbf{SDF} donne du
texte net à toute taille, au prix d'un atlas plus coûteux à produire.
\end{nkbilan}

\section{Exercices}

\begin{nkpratique}[1 — Le plus court chemin vers un texte à l'écran]
Dans une application NKGui minimale, remplacez le chargement de police par une
boucle sur les dix polices embarquées : pour chaque
\texttt{NkEmbeddedFontId} de 0 à 9, testez \texttt{IsAvailable}, récupérez son
nom avec \texttt{GetName}, et affichez la liste dans un panneau.

Puis ajoutez un bouton qui passe à la police suivante. Vous devrez recharger
l'atlas \emph{et} le ré-uploader dans le backend — l'exercice est là. Vérifiez ce
qui se passe si vous oubliez le ré-upload.
\end{nkpratique}

\begin{nkdemo}[2 — Voir l'atlas]
Récupérez le bitmap alpha8 avec \texttt{GetTexDataAsAlpha8}, enveloppez-le dans
une \texttt{NkImage} — le chapitre~5 vous a donné tout ce qu'il faut — et
sauvegardez-le en PNG.

Ouvrez le fichier. Vous verrez les glyphes empilés en étagères, avec deux pixels
de marge (\texttt{texGlyphPadding}). Comptez-les et comparez avec
\texttt{face->glyphCount}. Puis recommencez avec une plage de glyphes plus large
et observez l'atlas grandir. Enfin, mettez \texttt{sdfMode = true} et
regardez à quoi ressemble un champ de distance.
\end{nkdemo}

\begin{nkpratique}[3 — Écrire son propre rendu de texte]
Sans utiliser NKGui, écrivez une fonction qui prend une \texttt{NkFont*}, une
chaîne UTF-8, une position et une couleur, et remplit deux tableaux de sommets
et d'indices — exactement comme \texttt{NkGuiDrawList::AddText}. Testez-la sans
GPU : affichez pour chaque caractère son \emph{codepoint}, ses quatre coins et
ses UV.

Vérifiez ensuite que la somme des \texttt{advanceX} est bien égale à
\texttt{face->CalcTextSizeX(texte)}. Puis arrondissez volontairement
\texttt{advanceX} à l'entier et observez l'écart se creuser avec la longueur de
la chaîne : vous venez de reproduire le bug contre lequel le commentaire de
\nkfile{NkGuiDrawList.cpp:257-266} met en garde.
\end{nkpratique}

\begin{nkdemo}[4 — Toucher les plafonds]
Écrivez un programme console qui ajoute des polices à un atlas dans une boucle,
en incrémentant la taille d'un pixel à chaque tour, et qui s'arrête quand
\texttt{AddToAtlas} renvoie \texttt{nullptr}. Combien d'itérations ?

Ensuite, tentez d'ajouter une police avec
\texttt{GetGlyphRangesChineseFull()} et regardez ce que vaut
\texttt{face->glyphCount} après \texttt{Build()}. Comparez avec
\texttt{NK\_FONT\_MAX\_GLYPHS}. Que se passe-t-il pour les caractères au-delà ?
Cherchez la réponse dans le comportement observé, puis dans
\nkfile{NkFontAtlas.cpp} — et notez si le module vous a prévenu ou non.
\end{nkdemo}

\begin{nkquestions}
\begin{enumerate}
\item Quels sont les quatre temps de la séquence canonique ?
\item Que se passe-t-il si l'on demande un caractère après le \texttt{Build()} ?
\item Qu'est-ce qu'un atlas, et pourquoi n'affiche-t-on pas les glyphes un par
      un ?
\item Qu'apporte le mode SDF, et que coûte-t-il ?
\item Le dépôt embarque des polices : quel problème cela supprime-t-il ?
\end{enumerate}
\end{nkquestions}

\begin{nklogique}
Une application affiche correctement « Bonjour », mais les accents ressortent en
carrés vides. Le même texte s'affiche bien dans un autre écran de la même
application.

\begin{enumerate}
\item Énumérez au moins trois causes possibles.
\item Le fait que \emph{le même texte} s'affiche bien ailleurs élimine
      lesquelles ? Justifiez chaque élimination.
\item Décrivez la mesure qui tranche entre celles qui restent --- et dites ce
      qu'elle affiche dans chaque cas, sinon elle ne tranche rien.
\end{enumerate}
\end{nklogique}
